\documentclass[11pt,a4paper]{article}

% ---- Packages ----
\usepackage[utf8]{inputenc}
\usepackage[T1]{fontenc}
\usepackage{amsmath,amssymb,amsthm}
\usepackage{mathtools}
\usepackage[margin=1in,headheight=14pt]{geometry}
\usepackage{hyperref}
\usepackage{enumitem}
\usepackage{xcolor}
\usepackage{tcolorbox}
\usepackage{fancyhdr}
\usepackage{booktabs}

% ---- Hyperref ----
\hypersetup{
    colorlinks=true,
    linkcolor=red,
    citecolor=blue,
    urlcolor=blue
}

% ---- Header ----
\pagestyle{fancy}
\fancyhf{}
\fancyhead[L]{\textit{Lesson 12: Repeated Games and Incomplete Information}}
\fancyhead[R]{\thepage}
\renewcommand{\headrulewidth}{0.4pt}

% ---- Theorem Environments ----
\theoremstyle{definition}
\newtheorem{definition}{Definition}[section]
\newtheorem{example}{Example}[section]
\newtheorem{exercise}{Exercise}[section]

\theoremstyle{plain}
\newtheorem{theorem}{Theorem}[section]
\newtheorem{lemma}[theorem]{Lemma}
\newtheorem{proposition}[theorem]{Proposition}
\newtheorem{corollary}[theorem]{Corollary}

\theoremstyle{remark}
\newtheorem{remark}{Remark}[section]

% ---- Colored Boxes ----
\tcbuselibrary{theorems,skins,breakable}

\newtcolorbox{keyresult}[1][]{
    colback=green!5!white,
    colframe=green!50!black,
    fonttitle=\bfseries,
    title=#1,
    breakable
}

\newtcolorbox{rlconnection}[1][]{
    colback=blue!5!white,
    colframe=blue!50!black,
    fonttitle=\bfseries,
    title=#1,
    breakable
}

\newtcolorbox{warning}[1][]{
    colback=red!5!white,
    colframe=red!50!black,
    fonttitle=\bfseries,
    title=#1,
    breakable
}

% ---- Operators ----
\newcommand{\R}{\mathbb{R}}
\newcommand{\N}{\mathbb{N}}
\newcommand{\Q}{\mathbb{Q}}
\newcommand{\Z}{\mathbb{Z}}
\newcommand{\C}{\mathbb{C}}
\newcommand{\Prob}{\mathbb{P}}
\newcommand{\E}{\mathbb{E}}
\newcommand{\Var}{\mathrm{Var}}
\newcommand{\indicator}{\mathbf{1}}
\DeclareMathOperator{\supp}{supp}
\DeclareMathOperator{\conv}{conv}
\DeclareMathOperator{\argmax}{argmax}
\DeclareMathOperator{\argmin}{argmin}
\DeclareMathOperator{\BR}{BR}

% ---- Title ----
\title{Lesson 38: Signaling, Perfect Bayesian Equilibrium, and Potential Games}
\author{Curriculum: A Rigorous Learning Path to Multi-Agent Deep RL}
\date{}

\begin{document}
\maketitle
\tableofcontents
\newpage

\setcounter{section}{6}

\section{Perfect Bayesian Equilibrium}
\label{sec:pbe}

\subsection{Motivation}

In dynamic games of incomplete information, the BNE concept is too weak: it does not
constrain behavior at off-equilibrium information sets. We need a refinement that combines
sequential rationality with belief consistency.

\begin{definition}[System of Beliefs]
\label{def:beliefs}
In an extensive-form game with incomplete information, a \emph{system of beliefs}
$\mu$ assigns to each information set $h$ a probability distribution $\mu(\cdot | h)$
over the nodes in that information set.
\end{definition}

\begin{definition}[Perfect Bayesian Equilibrium]
\label{def:pbe}
A \emph{perfect Bayesian equilibrium} (PBE) is a pair $(\sigma, \mu)$ of a behavioral
strategy profile $\sigma$ and a belief system $\mu$ such that:
\begin{enumerate}[label=(\alph*)]
    \item \textbf{Sequential rationality:} At every information set $h$ of every player $i$,
    the strategy $\sigma_i$ is optimal given beliefs $\mu(\cdot | h)$ and the continuation
    strategies $\sigma_{-i}$.
    \item \textbf{Belief consistency:} At every information set $h$ that is reached with
    positive probability under $\sigma$, the beliefs $\mu(\cdot | h)$ are derived from
    $\sigma$ using Bayes' rule:
    \[
        \mu(x | h) = \frac{\Prob(x \text{ is reached} | \sigma)}
        {\sum_{x' \in h} \Prob(x' \text{ is reached} | \sigma)}
        \quad \text{for all } x \in h.
    \]
    \item At information sets reached with zero probability (off the equilibrium path),
    beliefs are unrestricted (any distribution over nodes in the information set is allowed).
\end{enumerate}
\end{definition}

\subsection{Signaling Games}

A particularly important class of dynamic Bayesian games is \emph{signaling games}.

\begin{definition}[Signaling Game]
\label{def:signaling-game}
A \emph{signaling game} has two players:
\begin{itemize}
    \item \textbf{Sender (S):} Has private type $\theta \in \Theta$, drawn from prior
    $p(\theta)$. Observes $\theta$, then sends a message $m \in M$.
    \item \textbf{Receiver (R):} Observes $m$ (but not $\theta$), then takes action
    $a \in A$.
\end{itemize}
Payoffs $u_S(\theta, m, a)$ and $u_R(\theta, m, a)$ depend on all three: type, message,
and action.
\end{definition}

\begin{definition}[Equilibrium Types in Signaling Games]
\label{def:equilibrium-types}
\begin{enumerate}[label=(\alph*)]
    \item A PBE is \emph{separating} if different types send different messages:
    $\sigma_S(\theta) \neq \sigma_S(\theta')$ for $\theta \neq \theta'$.
    \item A PBE is \emph{pooling} if all types send the same message:
    $\sigma_S(\theta) = m^*$ for all $\theta$.
    \item A PBE is \emph{semi-separating} (or partially pooling) if some types pool and
    others separate.
\end{enumerate}
\end{definition}

\begin{example}[Spence's Job Market Signaling]
\label{ex:spence}
A worker has ability $\theta \in \{L, H\}$ with $\Prob(\theta = H) = \pi$. The worker
chooses education level $e \geq 0$ (the ``signal''). An employer observes $e$ and offers
wage $w(e)$. Payoffs:
\begin{align*}
    u_S(\theta, e, w) &= w - c(\theta, e), \quad \text{with } c(H, e) < c(L, e)
    \text{ for } e > 0, \\
    u_R(\theta, e, w) &= \theta - w,
\end{align*}
where $c(\theta, e)$ is the cost of education (cheaper for high-ability workers:
\emph{single-crossing property}).

\textbf{Separating equilibrium.}
In a separating PBE, $e_H \neq e_L$ and the receiver's beliefs at $e_H$ are
$\mu(H | e_H) = 1$, at $e_L$ are $\mu(H | e_L) = 0$ (by Bayes' rule). The receiver offers
$w(e_H) = H$ and $w(e_L) = L$.

For this to be an equilibrium, each type must prefer its own signal:
\begin{align*}
    H - c(H, e_H) &\geq L - c(H, e_L) \quad \text{(type $H$ prefers $e_H$)}, \\
    L - c(L, e_L) &\geq H - c(L, e_H) \quad \text{(type $L$ prefers $e_L$)}.
\end{align*}
The single-crossing property ensures that such $e_H, e_L$ exist (typically $e_L = 0$ and
$e_H$ above a threshold).

\textbf{Pooling equilibrium.}
Both types choose $e^* = 0$. The receiver's belief at $e^*$ is the prior:
$\mu(H | e^*) = \pi$, so $w(e^*) = \pi H + (1-\pi) L$. Off-equilibrium beliefs at
$e \neq e^*$ can be chosen to deter deviation, e.g., $\mu(H | e) = 0$ for all $e \neq e^*$,
giving $w(e) = L$. Then no type wants to deviate if
$\pi H + (1-\pi)L - c(\theta, 0) \geq L - c(\theta, e)$ for all $e > 0$, which holds
when the pooling wage is sufficiently attractive.
\end{example}

\begin{proposition}[PBE Existence in Finite Signaling Games]
\label{prop:pbe-signaling-existence}
Every finite signaling game (finite $\Theta$, $M$, $A$) has a perfect Bayesian equilibrium.
\end{proposition}

\begin{proof}
A finite signaling game is a finite extensive-form game. The sequential equilibrium
existence theorem of Kreps and Wilson (1982) guarantees existence of a sequential
equilibrium, which is a refinement of PBE. In particular, every sequential equilibrium
is a PBE, so PBE exists.

Alternatively, one can construct a PBE directly: consider all possible sender strategies
$\sigma_S : \Theta \to \Delta(M)$. For each such strategy, compute the receiver's
beliefs via Bayes' rule at information sets reached with positive probability, and assign
arbitrary beliefs at unreached information sets. The receiver's optimal response defines
a best-response map. The sender's optimal message choice, given the receiver's response
function, defines another best-response map. By Kakutani's fixed point theorem (or Nash's
theorem applied to the agent normal form), a fixed point---hence a PBE---exists.
\end{proof}


%% ============================================================
%% SECTION 8: Potential Games
%% ============================================================

\section{Potential Games}
\label{sec:potential-games}

\subsection{Definitions}

Potential games form a special class of games where all players' incentives can be captured
by a single ``potential function.'' This structure has powerful implications for existence
of pure equilibria and convergence of learning dynamics.

\begin{definition}[Exact Potential Game]
\label{def:exact-potential}
A finite game $G = (N, (A_i)_{i \in N}, (u_i)_{i \in N})$ is an \emph{exact potential
game} if there exists a function $\Phi : A \to \R$ (the \emph{potential function}) such
that for every player $i$, every $a_{-i} \in A_{-i}$, and every $a_i, a_i' \in A_i$:
\[
    u_i(a_i, a_{-i}) - u_i(a_i', a_{-i}) = \Phi(a_i, a_{-i}) - \Phi(a_i', a_{-i}).
\]
\end{definition}

\begin{definition}[Ordinal Potential Game]
\label{def:ordinal-potential}
A finite game $G$ is an \emph{ordinal potential game} if there exists $\Phi : A \to \R$
such that for every player $i$, every $a_{-i}$, and every $a_i, a_i'$:
\[
    u_i(a_i, a_{-i}) > u_i(a_i', a_{-i}) \iff \Phi(a_i, a_{-i}) > \Phi(a_i', a_{-i}).
\]
\end{definition}

\begin{remark}
Every exact potential game is an ordinal potential game. The converse is false: ordinal
potentials only preserve the \emph{sign} of payoff differences, not their magnitude.
\end{remark}

\subsection{Existence of Pure Nash Equilibria}

\begin{theorem}[Potential Games Have Pure NE]
\label{thm:potential-pure-ne}
Every finite exact potential game (and every finite ordinal potential game) has at least one
pure Nash equilibrium.
\end{theorem}

\begin{proof}
Let $\Phi : A \to \R$ be the potential function. Since $A$ is finite, $\Phi$ attains its
maximum at some $a^* \in A$:
\[
    a^* \in \argmax_{a \in A} \Phi(a).
\]
We claim $a^*$ is a pure Nash equilibrium.

For any player $i$ and any deviation $a_i' \in A_i$:
\[
    \Phi(a_i^*, a_{-i}^*) \geq \Phi(a_i', a_{-i}^*)
\]
by the maximality of $a^*$. By the potential property:
\[
    u_i(a_i^*, a_{-i}^*) - u_i(a_i', a_{-i}^*)
    = \Phi(a_i^*, a_{-i}^*) - \Phi(a_i', a_{-i}^*) \geq 0.
\]
Thus $a_i^*$ is a best response to $a_{-i}^*$ for every player $i$, so $a^*$ is a NE.

For ordinal potential games, the same argument works: $\Phi(a_i^*, a_{-i}^*) \geq
\Phi(a_i', a_{-i}^*)$ implies $u_i(a_i^*, a_{-i}^*) \geq u_i(a_i', a_{-i}^*)$ by the
ordinal potential property (since if the strict inequality failed, i.e.,
$u_i(a_i', a_{-i}^*) > u_i(a_i^*, a_{-i}^*)$, then
$\Phi(a_i', a_{-i}^*) > \Phi(a_i^*, a_{-i}^*)$, contradicting maximality).
\end{proof}

\subsection{Convergence of Best-Response Dynamics}

\begin{definition}[Best-Response Dynamics]
\label{def:br-dynamics}
In \emph{(asynchronous) best-response dynamics}, players take turns updating their actions.
At each step, one player $i$ is selected (deterministically or randomly) and switches to a
best response to the current actions of all other players:
\[
    a_i^{t+1} \in \argmax_{a_i \in A_i} u_i(a_i, a_{-i}^t),
    \qquad a_j^{t+1} = a_j^t \text{ for } j \neq i.
\]
\end{definition}

\begin{theorem}[Finite Improvement Property]
\label{thm:fip}
In every finite exact potential game, best-response dynamics converge to a pure Nash
equilibrium in a finite number of steps.
\end{theorem}

\begin{proof}
We show that $\Phi$ strictly increases at each step where a player's action changes, hence
the dynamics must terminate (since $A$ is finite and $\Phi$ cannot increase indefinitely).

Suppose at step $t$, player $i$ switches from $a_i^t$ to $a_i^{t+1} \neq a_i^t$ with
$a_i^{t+1} \in \argmax_{a_i} u_i(a_i, a_{-i}^t)$. Since $a_i^{t+1}$ is a strictly
improving move (otherwise the player would not switch), we have:
\[
    u_i(a_i^{t+1}, a_{-i}^t) > u_i(a_i^t, a_{-i}^t).
\]
By the exact potential property:
\[
    \Phi(a_i^{t+1}, a_{-i}^t) - \Phi(a_i^t, a_{-i}^t)
    = u_i(a_i^{t+1}, a_{-i}^t) - u_i(a_i^t, a_{-i}^t) > 0.
\]

Thus $\Phi(a^{t+1}) > \Phi(a^t)$ whenever an action changes. Since $A$ is finite, the
sequence $\Phi(a^1), \Phi(a^2), \ldots$ is strictly increasing and takes values in the
finite set $\{\Phi(a) : a \in A\}$, so it must terminate. Termination occurs when no player
wants to change, i.e., at a pure Nash equilibrium.

\emph{Remark on step count:} The number of steps is bounded by $|A| = \prod_i |A_i|$
(since each action profile can appear at most once in the improvement path), but this
worst-case bound is typically very loose.
\end{proof}

\begin{corollary}
\label{cor:better-response}
The finite improvement property also holds for \emph{better-response dynamics} (where
player $i$ switches to any action that improves her payoff, not necessarily the best
response), since $\Phi$ still strictly increases.
\end{corollary}

\begin{proof}
If $u_i(a_i', a_{-i}^t) > u_i(a_i^t, a_{-i}^t)$, then
$\Phi(a_i', a_{-i}^t) > \Phi(a_i^t, a_{-i}^t)$ by the potential property. The same
termination argument applies.
\end{proof}

\subsection{Congestion Games}

The most important class of potential games in applications is congestion games.

\begin{definition}[Congestion Game]
\label{def:congestion-game}
A \emph{congestion game} is defined by $(N, E, (A_i)_{i \in N}, (d_e)_{e \in E})$ where:
\begin{itemize}
    \item $N = [n]$ is the player set.
    \item $E$ is a finite set of \emph{resources} (or \emph{facilities}).
    \item $A_i \subseteq 2^E$ is a collection of subsets of resources available to player
    $i$. Each action $a_i \in A_i$ is a set of resources that player $i$ uses.
    \item $d_e : \{1, \ldots, n\} \to \R$ is the \emph{delay function} (or cost function)
    for resource $e$, depending on the number of users.
\end{itemize}
The payoff (negative cost) to player $i$ under action profile $a = (a_1, \ldots, a_n)$ is:
\[
    u_i(a) = -\sum_{e \in a_i} d_e\bigl(n_e(a)\bigr),
\]
where $n_e(a) = |\{j \in N : e \in a_j\}|$ is the congestion (number of users) of
resource $e$.
\end{definition}

\begin{theorem}[Congestion Games are Exact Potential Games (Rosenthal)]
\label{thm:rosenthal}
Every congestion game is an exact potential game with potential function
\[
    \Phi(a) = -\sum_{e \in E} \sum_{k=1}^{n_e(a)} d_e(k).
\]
\end{theorem}

\begin{proof}
We must verify that for any player $i$, any $a_{-i}$, and any $a_i, a_i' \in A_i$:
\[
    u_i(a_i, a_{-i}) - u_i(a_i', a_{-i}) = \Phi(a_i, a_{-i}) - \Phi(a_i', a_{-i}).
\]

Consider the left side. Let $n_e = n_e(a_i, a_{-i})$ denote the congestion when player $i$
uses $a_i$, and $n_e' = n_e(a_i', a_{-i})$ when player $i$ uses $a_i'$. We partition the
resources into three groups:
\begin{enumerate}[label=(\roman*)]
    \item $e \in a_i \cap a_i'$: used by $i$ under both actions. Then $n_e = n_e'$, so
    $d_e(n_e)$ cancels in both $u_i$ and $\Phi$ differences.
    \item $e \in a_i \setminus a_i'$: used by $i$ under $a_i$ but not $a_i'$. Then
    $n_e = n_e' + 1$ (one more user under $a_i$). The contribution to $u_i(a_i, a_{-i})
    - u_i(a_i', a_{-i})$ is $-d_e(n_e) = -d_e(n_e' + 1)$. The contribution to
    $\Phi(a_i, a_{-i}) - \Phi(a_i', a_{-i})$ is $-\sum_{k=1}^{n_e'+1} d_e(k)
    + \sum_{k=1}^{n_e'} d_e(k) = -d_e(n_e'+1)$. These are equal.
    \item $e \in a_i' \setminus a_i$: used by $i$ under $a_i'$ but not $a_i$. Then
    $n_e' = n_e + 1$. The contribution to $u_i(a_i, a_{-i}) - u_i(a_i', a_{-i})$ is
    $0 - (-d_e(n_e')) = d_e(n_e + 1)$. The contribution to
    $\Phi(a_i, a_{-i}) - \Phi(a_i', a_{-i})$ is $-\sum_{k=1}^{n_e} d_e(k)
    + \sum_{k=1}^{n_e+1} d_e(k) = d_e(n_e+1)$. These are equal.
\end{enumerate}
Summing over all resources, $u_i(a_i, a_{-i}) - u_i(a_i', a_{-i})
= \Phi(a_i, a_{-i}) - \Phi(a_i', a_{-i})$.
\end{proof}

\begin{corollary}
\label{cor:congestion-ne}
Every finite congestion game has a pure Nash equilibrium, and best-response dynamics
converge to one in finitely many steps.
\end{corollary}

\begin{proof}
Immediate from Theorems~\ref{thm:rosenthal}, \ref{thm:potential-pure-ne},
and~\ref{thm:fip}.
\end{proof}

\begin{example}[Network Routing Game]
\label{ex:routing}
Consider a network with source $s$ and destination $d$, and $n$ users each choosing a path
from $s$ to $d$. Each edge $e$ has delay $d_e(k)$ depending on congestion $k$. This is a
congestion game: resources are edges, each player's action set is the set of $s$-$d$ paths,
and the total delay on a path is $\sum_{e \in \text{path}} d_e(n_e)$. By
Corollary~\ref{cor:congestion-ne}, a pure NE exists and can be found by best-response
dynamics.
\end{example}

\begin{example}[Resource Allocation]
\label{ex:resource-allocation}
Consider $n$ agents each choosing one of $m$ machines to process a job. Machine $e$ has
processing time $d_e(k) = k \cdot \ell_e$ (linear congestion) where $\ell_e$ is the base
load. Each agent's cost is the processing time on her chosen machine. This is a congestion
game with $A_i = E = \{1, \ldots, m\}$ (each action uses a single resource). The potential
function is $\Phi(a) = -\sum_{e=1}^m \ell_e \cdot \frac{n_e(n_e+1)}{2}$, and pure NE exist
with convergent best-response dynamics.
\end{example}

\begin{rlconnection}[Potential Games and Independent Learning in MARL]
Potential games are among the few game classes where independent learners (agents running
separate RL algorithms without coordination) are guaranteed to converge:
\begin{itemize}
    \item \textbf{Best-response dynamics} correspond to each agent independently computing
    optimal policies given fixed opponent policies. The finite improvement property ensures
    convergence.
    \item \textbf{Log-linear learning:} If each agent plays action $a_i$ with probability
    proportional to $\exp(\beta \cdot u_i(a_i, a_{-i}))$ (Boltzmann exploration), the
    joint process is a Markov chain that converges to the potential maximizer as
    $\beta \to \infty$.
    \item \textbf{Multi-agent reward shaping:} If a team reward function $\Phi$ can serve
    as a potential for individual reward functions $u_i$, then shaping individual rewards
    via $\Phi$ makes the game a potential game, guaranteeing convergence.
\end{itemize}
This is the theoretical basis for \emph{difference rewards} and \emph{Wonderful Life
Utility} in cooperative MARL.
\end{rlconnection}


%% ============================================================
%% SECTION 9: Connections to Multi-Agent RL
%% ============================================================

\section{Connections to Multi-Agent RL}
\label{sec:marl-connections}

We consolidate the connections between the game-theoretic results of this lesson and
modern multi-agent reinforcement learning.

\begin{rlconnection}[Independent Learners in Repeated Games]
When $n$ agents independently run Q-learning or policy gradient in a repeated game, each
agent treats the others as part of the environment. The folk theorems
(Theorems~\ref{thm:nash-folk} and~\ref{thm:spe-folk}) imply that the set of stable
outcomes is vast---any feasible, individually rational payoff can be a fixed point. This
explains the well-known instability and sensitivity to hyperparameters of independent
learners in general-sum games. Convergence guarantees for independent Q-learning exist
only in restrictive settings: two-player zero-sum games (minimax-Q), potential games, and
certain classes of team games.
\end{rlconnection}

\begin{rlconnection}[Self-Play and Fictitious Play]
Self-play---where an agent trains against copies of itself---is a form of fictitious play.
In zero-sum games, Robinson's theorem (Theorem~\ref{thm:robinson}) guarantees that the
empirical strategy profile converges to NE. Modern algorithms extend this:
\begin{itemize}
    \item \textbf{Neural Fictitious Self-Play (NFSP):} Uses neural networks to approximate
    best responses and average strategies, converging to approximate NE in large zero-sum
    games.
    \item \textbf{PSRO (Policy-Space Response Oracles):} Maintains a population of
    policies and computes a best response to the meta-mixture, directly implementing
    fictitious play in policy space.
    \item \textbf{Double Oracle:} Iteratively adds best responses to a restricted game,
    converging to the NE of the full game.
\end{itemize}
The regret bound of MWU (Theorem~\ref{thm:mwu-regret}) provides the convergence rate:
$O(1/\sqrt{T})$ for last-iterate convergence in zero-sum games with optimistic updates.
\end{rlconnection}

\begin{rlconnection}[Opponent Modeling as Bayesian Inference]
The Bayesian game framework (Section~\ref{sec:bayesian-games}) provides the foundation for
opponent modeling in MARL:
\begin{enumerate}[label=(\roman*)]
    \item \textbf{Type inference:} An agent maintains a belief $p(\theta_{-i} | h^t)$ over
    opponent types, updated via Bayes' rule as actions are observed.
    \item \textbf{Bayesian policy optimization:} The agent's policy maximizes expected
    return under the current belief: $\max_{\pi_i} \E_{\theta_{-i} \sim p(\cdot | h^t)}
    [R_i(\pi_i, \pi_{-i}(\theta_{-i}))]$.
    \item \textbf{Exploration-exploitation:} The agent may take actions to learn about
    opponents' types (information-gathering exploration), connecting to the theory of
    Bayesian exploration in MDPs.
\end{enumerate}
\end{rlconnection}

\begin{rlconnection}[Potential-Based Reward Shaping in Multi-Agent Settings]
Ng, Harada, and Russell's (1999) potential-based reward shaping theorem states that
adding a shaping reward $F(s, s') = \gamma \phi(s') - \phi(s)$ preserves the optimal
policy in single-agent MDPs. In multi-agent settings, this connects to potential games:
\begin{itemize}
    \item If the joint reward function serves as a potential for individual rewards,
    the multi-agent system is a potential game and independent learning converges.
    \item The \emph{difference reward} $D_i(a) = R(a) - R(a_{-i}, c_i)$ (where $c_i$ is a
    default action for agent $i$) is designed so that $D_i(a_i, a_{-i}) - D_i(a_i', a_{-i})
    = R(a_i, a_{-i}) - R(a_i', a_{-i})$, making $R$ an exact potential for the game with
    payoffs $D_i$.
    \item This theoretical property ensures that maximizing individual difference rewards
    aligns with maximizing the global objective, providing a principled approach to
    cooperative MARL.
\end{itemize}
\end{rlconnection}


%% ============================================================
%% SECTION 10: Exercises
%% ============================================================

\section{Exercises}
\label{sec:exercises}

\begin{exercise}[Finitely Repeated Prisoner's Dilemma]
\label{ex:finite-pd}
Consider the Prisoner's Dilemma from Example~\ref{ex:repeated-pd}, repeated $T = 100$
times. Show that the unique SPE is mutual defection in every period, regardless of $T$.
Explain why this result is ``paradoxical'' and how it motivates infinitely repeated games.
\end{exercise}

\begin{exercise}[Grim Trigger in a Different Game]
\label{ex:grim-trigger-general}
Consider the stage game:
\[
\begin{array}{c|ccc}
  & L & C & R \\ \hline
T & 6,6 & 0,8 & 0,0 \\
M & 8,0 & 4,4 & 0,0 \\
B & 0,0 & 0,0 & 2,2
\end{array}
\]
\begin{enumerate}[label=(\alph*)]
    \item Find all pure Nash equilibria of the stage game.
    \item Compute the minimax values $\underline{v}_1$ and $\underline{v}_2$.
    \item Show that the payoff $(6,6)$ (from $(T,L)$) can be sustained as an SPE of the
    infinitely repeated game using grim trigger with punishment $(B,R)$. Find the minimum
    discount factor.
    \item Can the payoff $(7,7)$ be sustained? If so, construct a strategy profile.
\end{enumerate}
\end{exercise}

\begin{exercise}[Tit-for-Tat Is Not SPE]
\label{ex:tft-not-spe}
Show that in the infinitely repeated Prisoner's Dilemma (Example~\ref{ex:repeated-pd}),
the tit-for-tat strategy profile is a Nash equilibrium for $\delta \geq 1/3$ but is
\emph{not} a subgame perfect equilibrium. (Hint: Consider the subgame after $(D, C)$.)
\end{exercise}

\begin{exercise}[Computing BNE]
\label{ex:compute-bne}
Consider a two-player Bayesian game where player 1 has type $\theta_1 \in \{H, L\}$ with
$\Prob(\theta_1 = H) = 1/2$, and player 2 has no private information. Actions:
$A_1 = A_2 = \{U, D\}$. Payoffs (for each type of player 1):
\[
\text{Type } H: \quad
\begin{array}{c|cc}
  & L & R \\ \hline
U & 2,1 & 0,0 \\
D & 0,0 & 1,2
\end{array}
\qquad \text{Type } L: \quad
\begin{array}{c|cc}
  & L & R \\ \hline
U & 0,0 & 1,2 \\
D & 2,1 & 0,0
\end{array}
\]
Find all Bayesian Nash equilibria (pure and mixed).
\end{exercise}

\begin{exercise}[Second-Price Auction as Bayesian Game]
\label{ex:second-price}
In a second-price sealed-bid auction with $n$ bidders having independent private values
$v_i \sim F$ (where $F$ is a continuous CDF on $[0, \bar{v}]$), show that bidding
$b_i = v_i$ (truthful bidding) is a dominant strategy, hence a BNE. Compare with the
first-price auction BNE from Example~\ref{ex:first-price-auction}.
\end{exercise}

\begin{exercise}[Verifying a Potential Function]
\label{ex:verify-potential}
Consider the two-player game with $A_1 = A_2 = \{X, Y\}$ and payoffs:
\[
\begin{array}{c|cc}
  & X & Y \\ \hline
X & 3,2 & 1,1 \\
Y & 2,1 & 4,3
\end{array}
\]
\begin{enumerate}[label=(\alph*)]
    \item Show this is an exact potential game by constructing the potential function $\Phi$.
    \item Find all pure Nash equilibria using $\Phi$.
    \item Verify that best-response dynamics converge from any initial profile.
\end{enumerate}
\end{exercise}

\begin{exercise}[MWU Convergence Rate]
\label{ex:mwu-rate}
Consider the zero-sum game Rock-Paper-Scissors:
\[
\begin{array}{c|ccc}
  & R & P & S \\ \hline
R & 0 & -1 & 1 \\
P & 1 & 0 & -1 \\
S & -1 & 1 & 0
\end{array}
\]
\begin{enumerate}[label=(\alph*)]
    \item What is the unique NE and the value of the game?
    \item If both players use MWU with $\eta = \sqrt{\ln 3 / T}$, bound the duality gap
    $\max_{a_1} g(a_1, \bar{\alpha}_2^T) - \min_{a_2} g(\bar{\alpha}_1^T, a_2)$ after
    $T = 1000$ rounds.
    \item How many rounds are needed for the duality gap to be below $0.01$?
\end{enumerate}
\end{exercise}

\begin{exercise}[Folk Theorem Application]
\label{ex:folk-application}
Consider the two-player stage game:
\[
\begin{array}{c|cc}
  & C & D \\ \hline
C & 5,5 & 0,7 \\
D & 7,0 & 2,2
\end{array}
\]
\begin{enumerate}[label=(\alph*)]
    \item Compute the feasible payoff set $V$ and the minimax values.
    \item Apply the Nash folk theorem to determine which payoff vectors are achievable as
    NE payoffs for large $\delta$.
    \item For the payoff $(5,5)$, find the minimum discount factor under grim trigger.
    \item For the payoff $(4,4)$, describe a strategy profile achieving it as an SPE and
    find the minimum $\delta$.
\end{enumerate}
\end{exercise}

\begin{exercise}[Congestion Game with Three Players]
\label{ex:congestion-three}
Three players share two resources $\{A, B\}$. Each player chooses exactly one resource.
Delay functions: $d_A(k) = k$ and $d_B(k) = 2k - 1$.
\begin{enumerate}[label=(\alph*)]
    \item Write the Rosenthal potential function $\Phi$.
    \item Enumerate all action profiles and compute $\Phi$ for each.
    \item Find all pure NE by maximizing $\Phi$ (equivalently, minimizing $-\Phi$).
    \item Verify that your NE are indeed equilibria by checking no player wants to deviate.
\end{enumerate}
\end{exercise}

\begin{exercise}[Separating vs.\ Pooling Equilibria]
\label{ex:signaling}
Consider a signaling game with two sender types $\theta \in \{G, B\}$ (good, bad),
$\Prob(\theta = G) = 1/3$, two messages $m \in \{H, L\}$ (high, low effort), and two
receiver actions $a \in \{A, R\}$ (accept, reject). Payoffs:
\begin{align*}
    u_S(G, H, A) &= 3, \quad u_S(G, H, R) = -1, \quad u_S(G, L, A) = 4, \quad
    u_S(G, L, R) = 0, \\
    u_S(B, H, A) &= 3, \quad u_S(B, H, R) = -2, \quad u_S(B, L, A) = 4, \quad
    u_S(B, L, R) = 1, \\
    u_R(\theta, m, A) &= \begin{cases} 1 & \text{if } \theta = G, \\
    -1 & \text{if } \theta = B, \end{cases} \quad
    u_R(\theta, m, R) = 0.
\end{align*}
\begin{enumerate}[label=(\alph*)]
    \item Does a separating PBE exist where type $G$ sends $H$ and type $B$ sends $L$?
    \item Does a pooling PBE exist where both types send $L$?
    \item Does a pooling PBE exist where both types send $H$?
\end{enumerate}
\end{exercise}


%% ============================================================
%% SECTION 11: Summary of Key Results
%% ============================================================

\section{Summary of Key Results}
\label{sec:summary}

\begin{center}
\renewcommand{\arraystretch}{1.4}
\begin{tabular}{@{}p{4.2cm}p{8cm}l@{}}
\toprule
\textbf{Result} & \textbf{Statement} & \textbf{Ref.} \\
\midrule

Unraveling Theorem &
Unique stage-game NE $\Rightarrow$ unique SPE of $G^T$: play NE every period &
Thm~\ref{thm:unraveling} \\

Nash Folk Theorem &
Any feasible, strictly individually rational payoff is a NE payoff for $\delta$ near 1 &
Thm~\ref{thm:nash-folk} \\

SPE Folk Theorem &
Same as Nash folk theorem but for SPE (under full dimensionality condition) &
Thm~\ref{thm:spe-folk} \\

Fictitious Play ($2 \times 2$) &
Empirical frequencies converge to NE in $2 \times 2$ zero-sum games &
Thm~\ref{thm:fp-2x2} \\

Robinson's Theorem &
Fictitious play converges to NE in all finite two-player zero-sum games &
Thm~\ref{thm:robinson} \\

MWU Regret Bound &
External regret $\leq 2\sqrt{T \ln |A_i|}$; hence MWU is no-regret &
Thm~\ref{thm:mwu-regret} \\

No-Regret $\Rightarrow$ NE &
If both players use no-regret, time-avg.\ play converges to NE in zero-sum &
Thm~\ref{thm:no-regret-ne} \\

BNE Existence &
Every finite Bayesian game has a BNE in mixed strategies &
Thm~\ref{thm:bne-existence} \\

PBE in Signaling Games &
Every finite signaling game has a PBE &
Prop~\ref{prop:pbe-signaling-existence} \\

Potential $\Rightarrow$ Pure NE &
Every finite potential game has a pure Nash equilibrium &
Thm~\ref{thm:potential-pure-ne} \\

Finite Improvement &
Best-response dynamics converge in finite steps in potential games &
Thm~\ref{thm:fip} \\

Rosenthal's Theorem &
Every congestion game is an exact potential game &
Thm~\ref{thm:rosenthal} \\

\bottomrule
\end{tabular}
\end{center}

\medskip

\begin{center}
\renewcommand{\arraystretch}{1.4}
\begin{tabular}{@{}p{3.5cm}p{8.5cm}@{}}
\toprule
\textbf{Game Class} & \textbf{MARL Implication} \\
\midrule

Repeated games (general-sum) &
Folk theorems $\Rightarrow$ many equilibria; independent learners may converge to any of them; sensitive to initialization and hyperparameters \\

Zero-sum games &
Minimax + no-regret $\Rightarrow$ convergent self-play (NFSP, PSRO, CFR) \\

Bayesian games &
Opponent modeling as Bayesian inference over types; principled exploration \\

Potential games &
Independent best-response / Q-learning converges; difference rewards make cooperative MARL a potential game \\

Congestion games &
Routing, resource allocation: pure NE exist, local search works \\

\bottomrule
\end{tabular}
\end{center}


\end{document}